\chapter{Хвильове рівняння: boundary та interior controllability як практичний міст}
\label{chap:wave-boundary-interior-bridge}

\sourcebox{Bridge-розділ. Це не повний source-preserved переклад розділів 2--3 Zuazua; відповідні \texttt{chapter02\_uk.tex} і \texttt{chapter03\_uk.tex} не були в завантаженому ZIP.}

Базова модель --- хвильове рівняння в обмеженій області $\Omega\subset\R^n$ з границею $\Gamma$:
\begin{equation}
\begin{cases}
 y_{tt}-\Delta y = 0, & (x,t)\in \Omega\times(0,T),\\
 y = 0, & (x,t)\in \Gamma\times(0,T),\\
 y(0)=y^0,\quad y_t(0)=y^1. &
\end{cases}
\end{equation}
Енергія природного розв'язку дорівнює
\begin{equation}
E(t)=\frac12\int_\Omega \left(|y_t|^2+|\nabla y|^2\right)\,\dd x.
\end{equation}
Для однорідної граничної умови Діріхле маємо збереження $E(t)=E(0)$. Це перший тест на правильність знаків, граничних умов і часової дискретизації.

\section{Граничне керування}

У boundary-control постановці керування діє на частині границі $\Gamma_0\subset\Gamma$:
\begin{equation}
\begin{cases}
 y_{tt}-\Delta y=0, & Q=\Omega\times(0,T),\\
 y=v, & \Sigma_0=\Gamma_0\times(0,T),\\
 y=0, & \Sigma_1=(\Gamma\setminus\Gamma_0)\times(0,T),\\
 y(0)=y^0,\quad y_t(0)=y^1. &
\end{cases}
\end{equation}
Точна керованість до нуля означає: для кожної пари початкових даних існує $v$ таке, що
\begin{equation}
 y(T)=0,\qquad y_t(T)=0.
\end{equation}
Через оборотність хвильового рівняння це еквівалентно можливості перейти між двома довільними допустимими станами за час $T$.

\section{HUM як операторна схема}

HUM-побудова замінює пошук керування на задачу спостережуваності для спряженої однорідної задачі. Нерівність спостережуваності має форму
\begin{equation}
 \|\phi(0)\|_{H^1_0(\Omega)}^2+\|\phi_t(0)\|_{L^2(\Omega)}^2
 \le C_T\int_0^T\int_{\Gamma_0}|\partial_\nu\phi|^2\,\dd\Gamma\,\dd t.
\end{equation}
Інженерна інтерпретація: щоб керувати всіма модами системи, дія на $\Gamma_0$ за час $T$ має бачити всю енергію спряженої хвилі.

\section{Внутрішнє локалізоване керування}

Для internal control керування діє в підобласті $\omega\subset\Omega$:
\begin{equation}
\begin{cases}
 y_{tt}-\Delta y = \chi_\omega u, & (x,t)\in Q,\\
 y=0, & (x,t)\in\Sigma,\\
 y(0)=y^0,\quad y_t(0)=y^1. &
\end{cases}
\end{equation}
Тоді спостережуваність для спряженої задачі має внутрішній вигляд:
\begin{equation}
 \|\phi(0)\|_{H^1_0(\Omega)}^2+\|\phi_t(0)\|_{L^2(\Omega)}^2
 \le C_T\int_0^T\int_\omega |\phi(x,t)|^2\,\dd x\,\dd t.
\end{equation}

\section{Практичні перевірки}

\begin{enumerate}[label=\arabic*.]
  \item Для некерованого рівняння енергія має зберігатися до чисельної похибки.
  \item Boundary control живе на $\Gamma_0\times(0,T)$, internal control --- на $\omega\times(0,T)$.
  \item Через скінченну швидкість поширення не можна вимагати керованості за довільно малий час.
  \item Невидимі моди або промені руйнують observability inequality.
\end{enumerate}
